% Part: second-order-logic
% Chapter: sol-and-sets
% Section: cardinalities

\documentclass[../../../include/open-logic-section]{subfiles}

\begin{document}

\olfileid{sol}{set}{crd}

\olsection{سیٹاں دی کارڈینیلٹی}

\begin{explain}
جیویں اسیں ایہ ظاہر کر سکدے آں کہ دائرہ متناہی یا لامتناہی اے،
!!{enumerable} یا !!{nonenumerable} اے، اویویں اسیں~$\Domain{M}$ دے کسے
ذیلی سیٹ دے متناہی یا لامتناہی، !!{enumerable} یا !!{nonenumerable} ہون دی خاصیت
دی تعریف وی کر سکدے آں۔
\end{explain}

\begin{prop}
ایہ فارمولا $\fn{Inf}(X) \ident$
\begin{multline*}
\lexists[u][(\lforall[x][\lforall[y][((X(x) \land X(y) \land \eq[u(x)][u(y)]) \lif
      \eq[x][y])]] \land \lforall[x][(X(x) \lif X(u(x)))] \land {} \\\lexists[y][(X(y) \land \lforall[x][(X(x)
      \lif \eq/[y][u(x)])])])]
\end{multline*}
متغیر قدر بندی~$s$ دے لحاظ نال اُدوں تے صرف اُدوں پورا اتردا اے جدوں
$s(X)$ لامتناہی ہووے۔
% PNB-CMP-078: X اُتے اک بہ اک خود نقشہ تے X وچ بندش شامل؛اصل global u نال متناہی X نوں وی لامتناہی آکھ سکدا سی۔
\end{prop}

\begin{prop}
ایہ فارمولا $\fn{Count}(X) \ident $
\begin{multline*}
(\lnot\lexists[x][X(x)] \lor {} \\
\lexists[z][\lexists[u][(X(z) \land
    \lforall[x][(X(x) \lif X(u(x)))] \land {} \\ \lforall[Y][((Y(z) \land
      \lforall[x][(Y(x) \lif Y(u(x)))]) \lif X = Y])]])
)
\end{multline*}
متغیر قدر بندی~$s$ دے لحاظ نال اُدوں تے صرف اُدوں پورا اتردا اے جدوں
$s(X)$، !!{enumerable} ہووے۔
% PNB-CMP-079: خالی سیٹ لئی الگ disjunct؛اصل z∈X غیر خالی enumerable سیٹ ہی گنّدا سی۔
\end{prop}

کانتور دے قضیے توں سانوں پتہ اے کہ !!{nonenumerable} سیٹ موجود نیں؛
اصل وچ لامتناہی سائز دے وکھ وکھ درجے نیں، تے اوہناں دی گنتی وی لامتناہی اے۔
سیٹ نظریہ سیٹاں دے سائز دا پورا حساب بناؤندا اے تے سیٹاں نوں
لامتناہی کارڈینیل اعداد دیندا اے۔ قدرتی اعداد متناہی سیٹاں دے سائز ماپن والے
کارڈینیل اعداد دا کم کردے نیں۔ !!{denumerable}
سیٹاں دی کارڈینیلٹی پہلی لامتناہی کارڈینیلٹی اے، جیہنوں
~$\aleph_0$ (”الیف-ناٹ“ یا ”الیف-زیرو“) آکھدے نیں۔ اگلا
لامتناہی سائز~$\aleph_1$ اے۔ ایہ اوہ سب توں چھوٹا سائز اے جیہڑا
کسے غیر قابلِ شمار سیٹ دا ہو سکدا اے؛ایہو جیہا سیٹ نہ متناہی ہوندا
اے نہ سائز~$\aleph_0$ رکھدا اے۔ اسیں ”$X$ دا سائز $\aleph_0$
اے“ ایویں تعریف کر سکدے آں:
$\fn{Aleph}_0(X) \liff \fn{Inf}(X) \land \fn{Count}(X)$۔
سیٹ~$X$ دا سائز~$\aleph_1$ اُدوں تے صرف اُدوں اے جدوں
اوہدا ہر اوہ ذیلی سیٹ جیہڑا سائز وچ اوہدے توں گھٹ اے، یا تاں
متناہی ہووے یا اوہدا سائز~$\aleph_0$ ہووے۔ اصل سیٹ آپ نہ
متناہی ہووے تے نہ اوہدا سائز~$\aleph_0$ ہووے۔
ایس لئی اسیں ایہ فارمولا لکھ سکدے آں:
$\fn{Aleph_1}(X) \ident \lforall[Y][((Y \subseteq X \land
  \lnot\cardeq{Y}{X}) \lif
  (\lnot\fn{Inf}(Y) \lor \fn{Aleph}_0(Y)))] \land \lnot
\fn{Count}(X)$۔ سائز~$\aleph_2$ وی ایس طرح تعریف ہوندا اے، وغیرہ۔
% PNB-CMP-080: Y=X نوں smaller-subset antecedent توں کڈھ کے cardinal strictness،تے X نوں nonenumerable منگیا؛اصل finite X قبول تے aleph_1 X رد کردا سی۔

اک سائز خاص دلچسپی رکھدا اے: استمرار (continuum) دی کارڈینیلٹی۔
ایہ~$\Pow{\Nat}$ دا سائز اے، یعنی اوہی سائز جیہڑا~$\Real$ دا اے۔
ایہ گل کہ کوئی سیٹ استمرار دے سائز دا اے، دوجے درجے دی منطق وچ
وی ظاہر کیتی جا سکدی اے، پر ایس لئی کجھ ہور کم کرنا پیندا اے۔

\end{document}
